% Options for packages loaded elsewhere
\PassOptionsToPackage{unicode}{hyperref}
\PassOptionsToPackage{hyphens}{url}
\documentclass[
  ignorenonframetext,
]{beamer}
\newif\ifbibliography
\usepackage{pgfpages}
\setbeamertemplate{caption}[numbered]
\setbeamertemplate{caption label separator}{: }
\setbeamercolor{caption name}{fg=normal text.fg}
\beamertemplatenavigationsymbolsempty
% remove section numbering
\setbeamertemplate{part page}{
  \centering
  \begin{beamercolorbox}[sep=16pt,center]{part title}
    \usebeamerfont{part title}\insertpart\par
  \end{beamercolorbox}
}
\setbeamertemplate{section page}{
  \centering
  \begin{beamercolorbox}[sep=12pt,center]{section title}
    \usebeamerfont{section title}\insertsection\par
  \end{beamercolorbox}
}
\setbeamertemplate{subsection page}{
  \centering
  \begin{beamercolorbox}[sep=8pt,center]{subsection title}
    \usebeamerfont{subsection title}\insertsubsection\par
  \end{beamercolorbox}
}
% Prevent slide breaks in the middle of a paragraph
\widowpenalties 1 10000
\raggedbottom
\AtBeginPart{
  \frame{\partpage}
}
\AtBeginSection{
  \ifbibliography
  \else
    \frame{\sectionpage}
  \fi
}
\AtBeginSubsection{
  \frame{\subsectionpage}
}
\usepackage{iftex}
\ifPDFTeX
  \usepackage[T1]{fontenc}
  \usepackage[utf8]{inputenc}
  \usepackage{textcomp} % provide euro and other symbols
\else % if luatex or xetex
  \usepackage{unicode-math} % this also loads fontspec
  \defaultfontfeatures{Scale=MatchLowercase}
  \defaultfontfeatures[\rmfamily]{Ligatures=TeX,Scale=1}
\fi
\usepackage{lmodern}
\ifPDFTeX\else
  % xetex/luatex font selection
\fi
% Use upquote if available, for straight quotes in verbatim environments
\IfFileExists{upquote.sty}{\usepackage{upquote}}{}
\IfFileExists{microtype.sty}{% use microtype if available
  \usepackage[]{microtype}
  \UseMicrotypeSet[protrusion]{basicmath} % disable protrusion for tt fonts
}{}
\makeatletter
\@ifundefined{KOMAClassName}{% if non-KOMA class
  \IfFileExists{parskip.sty}{%
    \usepackage{parskip}
  }{% else
    \setlength{\parindent}{0pt}
    \setlength{\parskip}{6pt plus 2pt minus 1pt}}
}{% if KOMA class
  \KOMAoptions{parskip=half}}
\makeatother
\usepackage{longtable,booktabs,array}
\newcounter{none} % for unnumbered tables
\usepackage{calc} % for calculating minipage widths
\usepackage{caption}
% Make caption package work with longtable
\makeatletter
\def\fnum@table{\tablename~\thetable}
\makeatother
\setlength{\emergencystretch}{3em} % prevent overfull lines
\providecommand{\tightlist}{%
  \setlength{\itemsep}{0pt}\setlength{\parskip}{0pt}}
\usepackage{bookmark}
\IfFileExists{xurl.sty}{\usepackage{xurl}}{} % add URL line breaks if available
\urlstyle{same}
\hypersetup{
  hidelinks,
  pdfcreator={LaTeX via pandoc}}

\author{\texorpdfstring{}{}}
\date{}

\begin{document}

\begin{frame}
\newpage
\end{frame}

\section{Common Pitfalls and What the Research
Shows}\label{sec:pitfalls}

\begin{frame}[allowframebreaks]{Common Pitfalls and What the Research
Shows}
The CIF research identifies recurring failure modes in multiagent
deployments. This section catalogs eight anti-patterns, each assessed
through what Parts 1 and 2 add to the problem and its mitigation.

These pitfalls are ranked by severity. We prioritize critical and
high-severity items as they represent verifiable vulnerabilities in the
defense architecture.
\end{frame}

\begin{frame}[allowframebreaks]{Pitfall 1: Implicit Trust (Critical)}
\protect\phantomsection\label{pitfall-1-implicit-trust-critical}
\textbf{The pattern}: Treating all inter-agent communication as trusted
by default.

\textbf{What the research shows}: Part 1's trust calculus exists
specifically because implicit trust enables trust laundering attacks.
Without explicit trust scores, a single compromised agent influences the
entire system---adversarial content enters through a low-trust source
and exits through a high-trust agent, with no mechanism to track or
attenuate the trust transfer.

Part 2's simulation confirms that trust exploitation attacks achieve
their highest success rates against architectures without trust decay.
The \(\delta^d\) mechanism isn't optional hardening---it's the
structural foundation that prevents systemic compromise from local
failures.

\textbf{Mitigation Strategies}:

\begin{enumerate}
\tightlist
\item
  Implement explicit trust scoring on every inter-agent channel
\item
  Require minimum trust thresholds for consequential actions
\item
  Apply delegation decay (\(\delta < 1\)) at every hop
\item
  Verify source identity on every inter-agent message
\end{enumerate}
\end{frame}

\begin{frame}[allowframebreaks]{Pitfall 2: Security as Afterthought
(Critical)}
\protect\phantomsection\label{pitfall-2-security-as-afterthought-critical}
\textbf{The pattern}: Adding cognitive security after the architecture
is finalized.

\textbf{What the research shows}: Part 1's defense composition algebra
demonstrates that security mechanisms compose best when designed
together. Retrofitted defenses create integration gaps---bypass
opportunities at the boundaries between the original architecture and
the bolted-on security layer.

Part 2's architecture-specific results show that systems with natural
security affordances (LangGraph's explicit state transitions, CrewAI's
role boundaries) outperform systems where security must be externally
imposed. Architecture is security posture.

\textbf{Mitigation Strategies}:

\begin{enumerate}
\tightlist
\item
  Define trust boundaries during architectural design, not deployment
\item
  Embed trust checks in delegation logic from the start
\item
  Build provenance tracking into belief management
\item
  Include cognitive security constraints in agent system prompts
\end{enumerate}
\end{frame}

\begin{frame}[allowframebreaks]{Pitfall 3: Uncalibrated Thresholds
(High)}
\protect\phantomsection\label{pitfall-3-uncalibrated-thresholds-high}
\textbf{The pattern}: Setting security thresholds without understanding
the tradeoffs.

\textbf{What the research shows}: Part 2's parametric simulation reveals
that detection rates are highly sensitive to threshold configuration.
The same defense module can range from too-strict (high false positive
rate, operational friction) to too-permissive (attacks succeed
undetected) depending on threshold settings.

The risk-profile-based configuration in Section 5 provides calibrated
starting points. But these are starting points---production thresholds
should be tuned against representative attack samples from Part 2's
corpus.

\textbf{Mitigation Strategies}:

\begin{enumerate}
\tightlist
\item
  Assess risk profile before configuring (Section 5)
\item
  Test thresholds against representative attack samples
\item
  Monitor false positive/negative rates in production
\item
  Adjust based on operational feedback
\end{enumerate}
\end{frame}

\begin{frame}[allowframebreaks]{Pitfall 4: Individual-Only Security
(High)}
\protect\phantomsection\label{pitfall-4-individual-only-security-high}
\textbf{The pattern}: Focusing on single-agent security while ignoring
multi-agent attack surfaces.

\textbf{What the research shows}: Part 1's entire contribution is
premised on the observation that multiagent systems introduce
qualitatively new attack surfaces. The adversary hierarchy
(\(\Omega_3\)--\(\Omega_5\)) specifically targets coordination,
consensus, and systemic properties that don't exist in single-agent
systems.

Part 2's results show that coordination attacks (sybil, timing, quorum
manipulation) are the \emph{hardest} to detect---precisely because they
exploit emergent properties of the agent collective rather than
vulnerabilities in individual agents.

\textbf{Mitigation Strategies}:

\begin{enumerate}
\tightlist
\item
  Implement Byzantine consensus for critical collective decisions
\item
  Require agent authentication before counting votes
\item
  Monitor for unusual coordination patterns (simultaneous votes,
  identical analyses)
\item
  Set quorum requirements assuming adversarial presence
\end{enumerate}
\end{frame}

\begin{frame}[allowframebreaks]{Pitfall 5: Static Tripwires (Medium)}
\protect\phantomsection\label{pitfall-5-static-tripwires-medium}
\textbf{The pattern}: Deploying canary tripwires once without rotation.

\textbf{What the research shows}: Part 1 notes that tripwire
effectiveness depends on unpredictability. If an adversary can identify
and avoid canary beliefs, the detection mechanism fails
silently---giving false confidence in detection coverage.

The analog in traditional security is static honeypots: effective
initially but useless once adversaries learn to recognize them.

\textbf{Mitigation Strategies}:

\begin{enumerate}
\tightlist
\item
  Implement automated canary rotation on a defined schedule
\item
  Vary placement across agents and belief categories
\item
  Monitor canary check patterns, not just modifications
\item
  Include non-obvious canaries that don't follow predictable naming
  conventions
\end{enumerate}
\end{frame}

\begin{frame}[allowframebreaks]{Pitfall 6: Ignoring Progressive Drift
(Medium)}
\protect\phantomsection\label{pitfall-6-ignoring-progressive-drift-medium}
\textbf{The pattern}: Only alerting on large, sudden belief changes.

\textbf{What the research shows}: Part 1's stealth-impact tradeoff
theorem bounds \emph{per-interaction} impact but explicitly acknowledges
that progressive drift---sub-threshold changes accumulating over
time---is the hardest attack pattern to detect. The theorem doesn't rule
out slow drift; it rules out sudden, invisible, high-impact attacks.

Part 2's multi-turn social engineering category, which achieves the
lowest detection rate (\textasciitilde73\%), partially exploits this
gap: attacks spread across multiple turns avoid the concentrated
statistical signature that single-turn attacks produce.

\textbf{Mitigation Strategies}:

\begin{enumerate}
\tightlist
\item
  Use sliding window drift detection (not just per-update thresholds)
\item
  Track \emph{cumulative} drift, not just per-update delta
\item
  Periodic baseline comparison over extended time windows
\item
  Alert on \emph{trends} as well as absolute magnitude
\end{enumerate}
\end{frame}

\begin{frame}[allowframebreaks]{Pitfall 7: Insufficient Logging
(Medium)}
\protect\phantomsection\label{pitfall-7-insufficient-logging-medium}
\textbf{The pattern}: Retaining insufficient information for
post-incident analysis.

\textbf{What the research shows}: Part 1's provenance tracking and
belief state representation are designed to produce a complete audit
trail. Without this trail, incident responders cannot reconstruct attack
paths, identify injection points, or assess the full scope of belief
corruption.

This is the cognitive security equivalent of running a production system
without structured logging. When something goes wrong---and it
will---the ability to investigate depends entirely on what was recorded.

\textbf{Mitigation Strategies}:

\begin{enumerate}
\tightlist
\item
  Log all belief updates with provenance tags (source, trust level,
  derivation chain)
\item
  Retain inter-agent message history
\item
  Take periodic cognitive state snapshots
\item
  Use structured logging formats that support causal analysis
\end{enumerate}
\end{frame}

\begin{frame}[allowframebreaks]{Pitfall 8: Single-Orchestrator Reliance
(High)}
\protect\phantomsection\label{pitfall-8-single-orchestrator-reliance-high}
\textbf{The pattern}: Relying entirely on one orchestrator's integrity
without backup.

\textbf{What the research shows}: Part 1's \(\Omega_5\) adversary
class---systemic compromise---is defined precisely as orchestrator-level
control. Section 4's Scenario 5 illustrates the consequences: the
attacker controls the entity responsible for coordinating defense,
rendering downstream defenses moot.

Part 2's hierarchical architecture results confirm the pattern:
hierarchical systems show strong \emph{average} detection because the
orchestrator is a natural defense chokepoint, but they have the worst
\emph{catastrophic} failure mode because that same chokepoint is a
single point of failure.

\textbf{Mitigation Strategies}:

\begin{enumerate}
\tightlist
\item
  Consider multi-orchestrator architectures for critical decisions
\item
  Monitor orchestrator behavior with the same rigor applied to worker
  agents
\item
  Workers should verify orchestrator identity on critical commands
\item
  Implement orchestrator-specific tripwires with rotation
\end{enumerate}
\end{frame}

\begin{frame}[allowframebreaks]{Summary Checklist}
\protect\phantomsection\label{summary-checklist}
{\def\LTcaptype{none} % do not increment counter
\begin{longtable}[]{@{}
  >{\raggedright\arraybackslash}p{(\linewidth - 4\tabcolsep) * \real{0.3333}}
  >{\raggedright\arraybackslash}p{(\linewidth - 4\tabcolsep) * \real{0.3333}}
  >{\raggedright\arraybackslash}p{(\linewidth - 4\tabcolsep) * \real{0.3333}}@{}}
\toprule\noalign{}
\begin{minipage}[b]{\linewidth}\raggedright
Pitfall
\end{minipage} & \begin{minipage}[b]{\linewidth}\raggedright
Assessment
\end{minipage} & \begin{minipage}[b]{\linewidth}\raggedright
Status
\end{minipage} \\
\midrule\noalign{}
\endhead
Implicit trust & Trust scoring implemented? & \(\square\) \\
Security afterthought & Security in initial architecture? &
\(\square\) \\
Uncalibrated thresholds & Thresholds tested against attacks? &
\(\square\) \\
Individual-only security & Byzantine consensus deployed? &
\(\square\) \\
Static tripwires & Canary rotation scheduled? & \(\square\) \\
Ignoring drift & Progressive drift monitoring? & \(\square\) \\
Insufficient logging & Full belief history retained? & \(\square\) \\
Single orchestrator & Orchestrator monitored? & \(\square\) \\
\bottomrule\noalign{}
\end{longtable}
}

Address unchecked items before production deployment.
\end{frame}

\end{document}
